\documentclass[11pt]{beamer}

% =========================================================================
%  Redondeo y ajuste por mínimos cuadrados ponderados
%  Plantilla: la de los seminarios del laboratorio (tema Warsaw, 4:3,
%  títulos en versalitas, bloques redondeados con sombra).
% =========================================================================

\usetheme{Warsaw}
\setbeamertemplate{blocks}[rounded][shadow=true]
\setbeamerfont{title}{family=\rmfamily, shape=\scshape}
\setbeamerfont{frametitle}{family=\rmfamily, shape=\scshape}
\setbeamercolor{block body}{bg=structure.fg!7!bg}
\setbeamertemplate{footline}{%
  \leavevmode%
  \hbox{%
  \begin{beamercolorbox}[wd=.5\paperwidth,ht=2.5ex,dp=1.125ex,right]{author in head/foot}%
    \usebeamerfont{author in head/foot}Redondeo y ajuste ponderado\hspace*{2ex}
  \end{beamercolorbox}%
  \begin{beamercolorbox}[wd=.5\paperwidth,ht=2.5ex,dp=1.125ex,right]{title in head/foot}%
    \usebeamerfont{title in head/foot}\insertframenumber\hspace*{2ex}
  \end{beamercolorbox}}%
  \vskip0pt}

% --- Idioma y tipografía ------------------------------------------------------
\usepackage[utf8]{inputenc}
\usepackage[T1]{fontenc}
\usepackage[spanish,es-tabla,es-nodecimaldot]{babel}
\usepackage{lmodern}
\usepackage[scaled=0.82]{beramono}
\usepackage{amsmath, amssymb}
\usepackage{siunitx}
\sisetup{output-decimal-marker={,}, group-minimum-digits=4, group-separator={\,},
         per-mode=symbol}
\usepackage{booktabs}

% --- Gráficos -------------------------------------------------------------------
\usepackage{tikz}
\usetikzlibrary{babel, arrows.meta, calc}
\usepackage{pgfplots}
\pgfplotsset{compat=1.17}
\definecolor{azulprof}{HTML}{1B365D}
\definecolor{azulacero}{HTML}{3B6EA8}
\definecolor{azulclaro}{HTML}{B7CBE3}
\definecolor{grispizarra}{HTML}{5B6770}
\definecolor{grisclaro}{HTML}{E4E8EC}
\definecolor{acento}{HTML}{B8800A}
% --- Estilo de gráficas de datos (como en «Leyes básicas de la electrostática»)
\pgfplotsset{
  datos/.style={
    tick label style={font=\scriptsize}, label style={font=\small},
    axis lines=left, axis line style={grispizarra},
    grid=major, grid style={grisclaro, line width=0.4pt},
    every tick/.style={grispizarra},
    legend style={font=\scriptsize, draw=none, fill=white, fill opacity=0.9,
                  text opacity=1, cells={anchor=west}},
    legend cell align=left,
    /pgf/number format/use comma, /pgf/number format/1000 sep={\,},
    scaled ticks=false,
  },
  puntoA/.style={only marks, mark=*, mark size=2.1pt, azulprof,
    mark options={fill=azulprof, draw=white, line width=0.3pt},
    error bars/.cd, x dir=both, x explicit, y dir=both, y explicit,
    error bar style={line width=0.6pt, azulprof}, error mark options={rotate=90, mark size=2.2pt, line width=0.6pt, azulprof}},
  puntoB/.style={only marks, mark=square*, mark size=1.9pt, azulacero,
    mark options={fill=azulclaro, draw=azulacero, line width=0.5pt},
    error bars/.cd, x dir=both, x explicit, y dir=both, y explicit,
    error bar style={line width=0.6pt, azulacero}, error mark options={rotate=90, mark size=2.2pt, line width=0.6pt, azulacero}},
  ajuste/.style={azulprof, line width=1.1pt},
  teorica/.style={acento, line width=0.9pt, densely dashed},
}

% Referencia a un panel de figura: «figura 5\,(b)», enlazada en bloque
\newcommand{\figpanel}[2]{\hyperref[#1]{figura~\ref*{#1}\,\textup{(#2)}}}



% Logo en la cabecera, como en la plantilla
\addtobeamertemplate{headline}{}{%
  \begin{tikzpicture}[remember picture, overlay]
    \node[anchor=north west, inner sep=0] at ($(current page.north west)+(0.18cm,-0.12cm)$)
      {\includegraphics[height=0.62cm]{logo_usal_negativo.png}};
  \end{tikzpicture}}

% --- Código: el mismo aspecto que en las explicaciones -------------------------
\usepackage{listings}
\usepackage[most]{tcolorbox}
\tcbuselibrary{listings}
\definecolor{codfondo}{HTML}{F7F7F8}
\definecolor{codborde}{HTML}{DADADF}
\definecolor{codtexto}{HTML}{383A42}
\definecolor{codclave}{HTML}{A626A4}   % palabras clave
\definecolor{codfunc}{HTML}{4078F2}    % funciones
\definecolor{codinterna}{HTML}{0184BC} % funciones internas
\definecolor{codcadena}{HTML}{50A14F}  % cadenas
\definecolor{codnumero}{HTML}{986801}  % números
\definecolor{codcoment}{HTML}{A0A1A7}  % comentarios
\definecolor{codinline}{HTML}{F2EFF1}
\definecolor{codrojo}{HTML}{C0365A}
\lstdefinestyle{python}{
  language=Python,
  basicstyle=\ttfamily\footnotesize\color{codtexto},
  keywordstyle=\color{codclave},
  keywordstyle=[2]\color{codinterna},
  keywordstyle=[3]\color{codfunc},
  morekeywords=[1]{lambda, None, True, False},
  morekeywords=[2]{abs, int, float, max, round, len, range, sum, dict},
  morekeywords=[3]{redondea, coma, ajuste_ponderado, floor, log10, polyfit, sqrt,
                   asarray, ppf, replace},
  stringstyle=\color{codcadena},
  commentstyle=\color{codcoment}\itshape,
  showstringspaces=false, upquote=true,
  columns=fixed, basewidth=0.5em, keepspaces=true,
  breaklines=true, breakindent=2em,
  literate=
    {á}{{\'a}}1 {é}{{\'e}}1 {í}{{\'i}}1 {ó}{{\'o}}1 {ú}{{\'u}}1 {ñ}{{\~n}}1
    {0}{{\textcolor{codnumero}{0}}}1 {1}{{\textcolor{codnumero}{1}}}1
    {2}{{\textcolor{codnumero}{2}}}1 {3}{{\textcolor{codnumero}{3}}}1
    {4}{{\textcolor{codnumero}{4}}}1 {5}{{\textcolor{codnumero}{5}}}1
    {6}{{\textcolor{codnumero}{6}}}1 {7}{{\textcolor{codnumero}{7}}}1
    {8}{{\textcolor{codnumero}{8}}}1 {9}{{\textcolor{codnumero}{9}}}1,
}
\newtcblisting{codigo}[1][]{%
  listing only, enhanced, colback=codfondo, colframe=codborde, boxrule=0.5pt,
  arc=4pt, left=6pt, right=6pt, top=2pt, bottom=2pt,
  before skip=4pt, after skip=6pt,
  listing options={style=python}, #1}
\newcommand{\cod}[1]{{\setlength{\fboxsep}{1.4pt}\colorbox{codinline}{\textcolor{codrojo}{\ttfamily\small #1}}}}

% Texto resaltado en azul, como los términos clave de la plantilla
\newcommand{\resalta}[1]{\textcolor{structure.fg!75!white}{#1}}

\newcommand{\vs}{v_{\mathrm{s}}}
\newcommand{\vst}{\tilde{v}_{\mathrm{s}}}
\newcommand{\fE}{f_{\mathrm{E}}}
\newcommand{\fR}{f_{\mathrm{R}}}

\title[Redondeo y ajuste ponderado]{Redondeo y ajuste por mínimos cuadrados ponderados}
\author{Daniel Martín Guerrero}
\institute{Laboratorio de Mecánica y Ondas, Grado en Física\\Universidad de Salamanca}
\date{octubre de 2026}
\titlegraphic{%
  \begin{tikzpicture}
    \draw[structure.fg!35!white, line width=0.5pt] (0,0) -- (8,0);
    \draw[structure.fg!80!white, line width=0.9pt, domain=0:8, samples=320, smooth]
      plot (\x, {0.22*sin(deg(\x*(2.4+0.55*\x)))});
  \end{tikzpicture}}

% =========================================================================
\begin{document}

\begin{frame}
  \titlepage
\end{frame}

\begin{frame}{Índice}
  \tableofcontents
\end{frame}

% =========================================================================
\section{Objetivo}
% =========================================================================
\begin{frame}{Objetivo}
  Cuando hacemos una gráfica con ajuste, el código se encarga de dos tareas:
  \begin{block}{}
    \begin{itemize}
      \item \resalta{Calcular la recta} que mejor describe los datos, teniendo en
        cuenta lo fiable que es cada punto.
      \item \resalta{Presentar el resultado} con el número correcto de cifras
        significativas.
    \end{itemize}
  \end{block}
  Vamos a recorrer el código línea a línea con los números reales del ajuste de
  la velocidad del sonido (caso~1 del informe): $v$ frente a $x=1-\fE/\fR$.
\end{frame}

% =========================================================================
\section{El redondeo}
% =========================================================================
\begin{frame}{Criterio de redondeo}
  \begin{block}{}
    \resalta{Criterio del informe:}
    \begin{itemize}
      \item la incertidumbre se da con \emph{una} cifra significativa;
      \item si esa cifra es un 1, se conservan \emph{dos};
      \item el valor se redondea a la misma posición decimal.
    \end{itemize}
  \end{block}
  \centering\small
  \begin{tabular}{@{}lcl@{}}
    \toprule
    Sin redondear & 1.ª cifra & Resultado\\
    \midrule
    A: $0{,}353\,17\pm0{,}011\,8$ & 1 & $0{,}353\pm0{,}012$\\
    B: $244{,}66\pm88{,}6$ & 8 & $(2{,}4\pm0{,}9)\times10^{2}$\\
    \bottomrule
  \end{tabular}
  \par\vspace{4pt}
  \raggedright\small Los ejemplos A y B nos acompañarán en toda esta sección.
\end{frame}

\begin{frame}[fragile]{La función \texttt{redondea}}
  \begin{codigo}[listing options={style=python, basicstyle=\ttfamily\scriptsize\color{codtexto}}]
def redondea(valor, incertidumbre):
    exp = math.floor(math.log10(abs(incertidumbre)))
    primera = int(abs(incertidumbre) / 10**exp)
    decimales = -(exp - 1) if primera == 1 else -exp
    v, u = round(valor, decimales), round(incertidumbre, decimales)
    d = max(decimales, 0)
    return f"{v:.{d}f}".replace(".", ","), f"{u:.{d}f}".replace(".", ",")
  \end{codigo}
  Recibe el valor y su incertidumbre \resalta{sin redondear} y devuelve dos
  textos listos para el informe. Los pasos son:
  \begin{enumerate}
    \item localizar la primera cifra significativa;
    \item ver si es un 1;
    \item decidir dónde cortar;
    \item redondear y escribir con coma decimal.
  \end{enumerate}
\end{frame}

\begin{frame}[fragile]{Paso 1: posición de la primera cifra}
  \begin{codigo}
exp = math.floor(math.log10(abs(incertidumbre)))
  \end{codigo}
  El \resalta{logaritmo decimal} responde a la pregunta: \emph{¿a qué número hay
  que elevar 10 para obtener este?}
  \begin{columns}[T]
    \begin{column}{0.4\textwidth}
      \centering\small
      \begin{tabular}{@{}rcc@{}}
        \toprule
        Número & Potencia & $\log_{10}$\\
        \midrule
        100 & $10^{2}$ & 2\\
        10 & $10^{1}$ & 1\\
        1 & $10^{0}$ & 0\\
        0,1 & $10^{-1}$ & $-1$\\
        0,01 & $10^{-2}$ & $-2$\\
        \bottomrule
      \end{tabular}
    \end{column}
    \begin{column}{0.56\textwidth}
      \small Si el número cae entre dos potencias, su logaritmo cae entre los dos
      enteros:
      \begin{itemize}
        \item $\log_{10}(88{,}6)\approx 1{,}95$;
        \item $\log_{10}(0{,}0118)\approx -1{,}93$.
      \end{itemize}
      \cod{abs(...)} garantiza que el número sea positivo.
    \end{column}
  \end{columns}
\end{frame}

\begin{frame}{Paso 1: el papel de \texttt{floor}}
  \cod{math.floor} redondea \resalta{siempre hacia abajo} y conserva la parte
  entera, que indica la posición de la primera cifra significativa:
  \begin{block}{}
    \centering
    \begin{tabular}{@{}lccl@{}}
      & $\log_{10}$ & \texttt{floor} & Posición\\
      \midrule
      A: $0{,}0118$ & $-1{,}93$ & $-2$ & centésimas ($10^{-2}$)\\
      B: $88{,}6$ & $1{,}95$ & $1$ & decenas ($10^{1}$)\\
    \end{tabular}
  \end{block}
  \begin{itemize}
    \item En los negativos, «hacia abajo» significa \emph{alejándose del cero}:
      $-1{,}93$ pasa a $-2$, no a $-1$, y así debe ser, porque el 1 de $0{,}0118$
      está en las centésimas.
    \item A ojo es fácil ver dónde empieza un número; el logaritmo es la forma
      matemática de que el ordenador lo averigüe para cualquier valor.
  \end{itemize}
\end{frame}

\begin{frame}[fragile]{Paso 2: valor de la primera cifra}
  \begin{codigo}
primera = int(abs(incertidumbre) / 10**exp)
  \end{codigo}
  Dividir entre $10^{\text{exp}}$ desplaza la coma hasta dejar la primera cifra
  delante, e \cod{int(...)} corta todo lo que hay detrás:
  \begin{block}{}
    \begin{equation}
      \text{A: }\frac{0{,}0118}{10^{-2}}=1{,}18\;\rightarrow\;1,
      \qquad
      \text{B: }\frac{88{,}6}{10^{1}}=8{,}86\;\rightarrow\;8.
    \end{equation}
  \end{block}
  Ya sabemos \resalta{dónde está} la primera cifra (\texttt{exp}) y
  \resalta{cuál es} (\texttt{primera}): con eso basta para aplicar el criterio.
\end{frame}

\begin{frame}[fragile]{Paso 3: posición de corte}
  \begin{codigo}
decimales = -(exp - 1) if primera == 1 else -exp
  \end{codigo}
  \begin{itemize}
    \item Si la primera cifra es un 1, se conservan \resalta{dos} cifras:
      se corta una posición más a la derecha, \cod{-(exp - 1)}.
    \item En otro caso se conserva \resalta{una}: \cod{-exp}.
  \end{itemize}
  \begin{block}{}
    \centering
    \begin{tabular}{@{}lccc@{}}
      & \texttt{exp} & \texttt{primera} & \texttt{decimales}\\
      \midrule
      A & $-2$ & 1 & $-(-2-1)=3$ \quad (milésimas)\\
      B & $1$ & 8 & $-1$ \quad (decenas)\\
    \end{tabular}
  \end{block}
  Un número de decimales negativo indica que se redondea a la izquierda de la
  coma.
\end{frame}

\begin{frame}[fragile]{Paso 4: redondeo y escritura}
  \begin{codigo}[listing options={style=python, basicstyle=\ttfamily\scriptsize\color{codtexto}}]
v, u = round(valor, decimales), round(incertidumbre, decimales)
d = max(decimales, 0)
return f"{v:.{d}f}".replace(".", ","), f"{u:.{d}f}".replace(".", ",")
  \end{codigo}
  \begin{itemize}
    \item \cod{round} redondea valor e incertidumbre a la \resalta{misma
      posición}.
    \item \cod{max(decimales, 0)}: si se redondea a las decenas, no hay
      decimales que escribir.
    \item El formato escribe exactamente \texttt{d} decimales y \cod{replace}
      cambia el punto por la coma.
  \end{itemize}
  \begin{block}{}
    \centering
    \begin{tabular}{@{}lll@{}}
      & Entrada & Devuelve\\
      \midrule
      A & $0{,}353\,17$; $0{,}011\,8$ & \texttt{0,353} y \texttt{0,012}\\
      B & $244{,}66$; $88{,}6$ & \texttt{240} y \texttt{90}, es decir,
      $(2{,}4\pm0{,}9)\times10^{2}$\\
    \end{tabular}
  \end{block}
\end{frame}

\begin{frame}[fragile]{La función \texttt{coma} y dos detalles}
  \begin{codigo}
def coma(x, d):
    return f"{x:.{d}f}".replace(".", ",")
  \end{codigo}
  Solo da formato, sin aplicar el criterio de redondeo; se usa para números
  auxiliares como $t_{95,\nu}$ o $\chi^{2}/\nu$.
  \begin{block}{}
    \resalta{Empates:} Python lleva los 5 exactos a la cifra par
    (\cod{round(0.125, 2)} da $0{,}12$), como indica la regla de las
    diapositivas de teoría de errores.
  \end{block}
  \begin{block}{}
    \resalta{Caso límite:} $u=0{,}0196$ empieza por 1 y conserva dos cifras,
    $0{,}020$; tras redondear empieza por 2, pero mantiene las dos. Es poco
    frecuente y, si aparece, se corrige a mano.
  \end{block}
\end{frame}

% =========================================================================
\section{El ajuste ponderado}
% =========================================================================
\begin{frame}{Motivación}
  Cada punto es la media de seis lecturas, con $u^{2}=s^{2}/n+\Delta^{2}/12$, y la
  abscisa $x=1-\fE/\fR$ hereda la incertidumbre de las frecuencias:
  \begin{center}\small
  \begin{tabular}{@{}lcc@{}}
    \toprule
    Serie & $u(y)$ (m/s) & $a\,u(x)$ (m/s)\\
    \midrule
    Máxima, acercamiento & 0,0026 & 0,030\\
    Máxima, alejamiento & 0,0012 & \resalta{0,184}\\
    Media, acercamiento & 0,0009 & 0,011\\
    Media, alejamiento & 0,0008 & 0,015\\
    Mínima, acercamiento & 0,0004 & 0,025\\
    Mínima, alejamiento & 0,0007 & 0,008\\
    \bottomrule
  \end{tabular}
  \end{center}
  \begin{block}{}
    La incertidumbre horizontal, trasladada a la vertical, es entre 10 y 150 veces
    mayor: el ajuste no puede ignorar la $x$ ni tratar igual todos los puntos.
  \end{block}
\end{frame}

\begin{frame}[fragile]{La función \texttt{ajuste\_ponderado} (I)}
  \begin{codigo}[listing options={style=python, basicstyle=\ttfamily\scriptsize\color{codtexto}}]
def ajuste_ponderado(x, ux, y, uy, modelo="libre"):
    x, ux, y, uy = (np.asarray(t, float) for t in (x, ux, y, uy))
    n = len(x)
    a = np.polyfit(x, y, 1)[0] if modelo == "libre" else (x @ y) / (x @ x)
    for _ in range(500):
        w = 1 / (uy**2 + a**2 * ux**2)
        if modelo == "libre":
            W = w.sum()
            m = lambda z: (w * z).sum() / W
            a_nueva = (m(x * y) - m(x) * m(y)) / (m(x * x) - m(x) ** 2)
        else:
            a_nueva = (w * x * y).sum() / (w * x * x).sum()
        if abs(a_nueva - a) <= 1e-12 * max(1.0, abs(a)):
            a = a_nueva
            break
        a = a_nueva
  \end{codigo}
\end{frame}

\begin{frame}[fragile]{La función \texttt{ajuste\_ponderado} (II)}
  \begin{codigo}[listing options={style=python, basicstyle=\ttfamily\scriptsize\color{codtexto}}]
ue = np.sqrt(uy**2 + a**2 * ux**2)
w = 1 / ue**2
if modelo == "libre":
    W = w.sum()
    m = lambda z: (w * z).sum() / W
    b = m(y) - a * m(x)
    ua = 1 / math.sqrt(W) / math.sqrt(m(x * x) - m(x) ** 2)
    ub = ua * math.sqrt(m(x * x))
    nu = n - 2
else:
    b, ub, nu = 0.0, 0.0, n - 1
    ua = 1 / math.sqrt((w * x * x).sum())
chi2 = float((w * (y - a * x - b) ** 2).sum())
factor = max(1.0, math.sqrt(chi2 / nu)) if nu > 0 else 1.0
t95 = float(stats.t.ppf(0.975, nu)) if nu > 0 else float("nan")
  \end{codigo}
\end{frame}

\begin{frame}[fragile]{Preparación y pendiente inicial}
  \begin{codigo}[listing options={style=python, basicstyle=\ttfamily\scriptsize\color{codtexto}}]
x, ux, y, uy = (np.asarray(t, float) for t in (x, ux, y, uy))
n = len(x)
a = np.polyfit(x, y, 1)[0] if modelo == "libre" else (x @ y) / (x @ x)
  \end{codigo}
  \begin{itemize}
    \item \cod{np.asarray} convierte las listas en \emph{arrays}, con los que se
      opera punto a punto: \cod{x * y} multiplica cada $x_i$ por su $y_i$.
    \item \cod{n = len(x)} cuenta los puntos: $n=6$.
    \item La iteración necesita una pendiente de partida:
      \begin{itemize}
        \item modelo libre: el ajuste sin pesos de \cod{np.polyfit};
        \item recta por el origen: $a=\sum x_iy_i\big/\sum x_i^{2}$.
      \end{itemize}
  \end{itemize}
  \begin{block}{}
    En el caso~1 la pendiente de partida es $a_{0}=82{,}4$~m/s: el ajuste básico,
    arrastrado por el punto anómalo.
  \end{block}
\end{frame}

\begin{frame}[fragile]{Pesos y varianza efectiva}
  \begin{codigo}
w = 1 / (uy**2 + a**2 * ux**2)
  \end{codigo}
  \begin{block}{}
    \resalta{Peso de cada punto:}
    \begin{equation}
      w_i=\frac{1}{u_i^{2}},\qquad u_i^{2}=u^{2}(y_i)+a^{2}\,u^{2}(x_i).
    \end{equation}
  \end{block}
  \begin{itemize}
    \item Un punto con el doble de incertidumbre pesa \resalta{cuatro veces
      menos}: los puntos precisos mandan en la recta.
    \item ¿Por qué $a^{2}$? Un desplazamiento horizontal $\delta x$ equivale, sobre
      una recta de pendiente $a$, a uno vertical $a\,\delta x$; sumando ambas
      contribuciones como varianzas independientes se obtiene $u_i^{2}$.
    \item Para calcular $w_i$ hace falta $a$, que es justo lo que se busca: por eso
      hay que \resalta{iterar}.
  \end{itemize}
\end{frame}

\begin{frame}[fragile]{La media ponderada}
  \begin{codigo}
W = w.sum()
m = lambda z: (w * z).sum() / W
  \end{codigo}
  \cod{lambda} define una función en una línea; \texttt{z} es solo el nombre de
  su argumento. \cod{m} calcula la \resalta{media ponderada}:
  \begin{block}{}
    \begin{equation}
      \langle z\rangle=\frac{\sum_i w_i\,z_i}{\sum_i w_i}.
    \end{equation}
  \end{block}
  \begin{columns}[T]
    \begin{column}{0.5\textwidth}
      \centering\small
      \begin{tabular}{@{}llr@{}}
        \toprule
        Llamada & Calcula & Caso 1\\
        \midrule
        \cod{m(x)} & $\langle x\rangle$ & $3{,}27\times10^{-4}$\\
        \cod{m(y)} & $\langle y\rangle$ & $-0{,}0214$\\
        \cod{m(x*x)} & $\langle x^{2}\rangle$ & $6{,}49\times10^{-7}$\\
        \cod{m(x*y)} & $\langle xy\rangle$ & $1{,}26\times10^{-4}$\\
        \bottomrule
      \end{tabular}
    \end{column}
    \begin{column}{0.46\textwidth}
      \small Ejemplo: con pesos $(1,1,2)$ y valores $(10,20,40)$,
      $\langle z\rangle=110/4=27{,}5$, mayor que la media aritmética
      ($23{,}3$), porque el 40 pesa el doble.
    \end{column}
  \end{columns}
\end{frame}

\begin{frame}[fragile]{La pendiente ponderada}
  \begin{codigo}[listing options={style=python, basicstyle=\ttfamily\scriptsize\color{codtexto}}]
a_nueva = (m(x * y) - m(x) * m(y)) / (m(x * x) - m(x) ** 2)
  \end{codigo}
  \begin{block}{}
    \resalta{Mínimos cuadrados ponderados:} minimizando
    $\chi^{2}=\sum_i w_i\,(y_i-ax_i-b)^{2}$ respecto de $a$ y $b$,
    \begin{equation}
      a=\frac{\langle xy\rangle-\langle x\rangle\langle y\rangle}
             {\langle x^{2}\rangle-\langle x\rangle^{2}},
      \qquad
      b=\langle y\rangle-a\langle x\rangle .
    \end{equation}
  \end{block}
  \begin{itemize}
    \item El numerador es una «covarianza» ponderada entre $x$ e $y$, y el
      denominador, una «varianza» ponderada de $x$.
    \item Con todos los pesos iguales se recuperan las fórmulas básicas de las
      notas de incertidumbres.
    \item Para la recta por el origen: \cod{a\_nueva = (w*x*y).sum() / (w*x*x).sum()}.
  \end{itemize}
\end{frame}

\begin{frame}[fragile]{Criterio de convergencia}
  \begin{codigo}
if abs(a_nueva - a) <= 1e-12 * max(1.0, abs(a)):
    a = a_nueva
    break
a = a_nueva
  \end{codigo}
  \begin{columns}[T]
    \begin{column}{0.52\textwidth}
      \small Si la pendiente nueva apenas difiere de la anterior (una billonésima
      relativa), el cálculo ha convergido y se sale con \cod{break}; si no, se
      guarda y se da otra vuelta. \cod{range(500)} es solo un tope de
      seguridad.
    \end{column}
    \begin{column}{0.44\textwidth}
      \centering\small
      \begin{tabular}{@{}lc@{}}
        \toprule
        Vuelta & $a$ (m/s)\\
        \midrule
        Inicio & 82,4\\
        1 & 244,279\\
        2 & 244,657\\
        3 & 244,657\\
        \bottomrule
      \end{tabular}
    \end{column}
  \end{columns}
  \begin{block}{}
    El salto de 82 a 244~m/s en la primera vuelta es el efecto de los pesos: el
    punto anómalo deja de arrastrar la recta.
  \end{block}
\end{frame}

\begin{frame}[fragile]{Pesos finales}
  \begin{codigo}
ue = np.sqrt(uy**2 + a**2 * ux**2)
w = 1 / ue**2
  \end{codigo}
  \begin{columns}[T]
  \begin{column}{0.58\textwidth}
  \centering\small
  \begin{tabular}{@{}lcr@{}}
    \toprule
    Serie & $u_i$ (m/s) & Peso\\
    \midrule
    Mínima, alejamiento & 0,008 & 51\,\%\\
    Media, acercamiento & 0,011 & 26\,\%\\
    Media, alejamiento & 0,015 & 14\,\%\\
    Mínima, acercamiento & 0,025 & 5\,\%\\
    Máxima, acercamiento & 0,030 & 3,5\,\%\\
    Máxima, alejamiento & 0,184 & \resalta{0,1\,\%}\\
    \bottomrule
  \end{tabular}
  \end{column}
  \begin{column}{0.38\textwidth}
  \small Con la pendiente ya convergida se recalculan la incertidumbre efectiva
  de cada punto y su peso definitivo.
  \begin{block}{}
    \small El punto anómalo prácticamente no cuenta; las series a velocidad media
    y mínima, las más precisas, deciden la recta.
  \end{block}
  \end{column}
  \end{columns}
\end{frame}

\begin{frame}[fragile]{Ordenada e incertidumbres de los parámetros}
  \begin{codigo}[listing options={style=python, basicstyle=\ttfamily\scriptsize\color{codtexto}}]
b = m(y) - a * m(x)
ua = 1 / math.sqrt(W) / math.sqrt(m(x * x) - m(x) ** 2)
ub = ua * math.sqrt(m(x * x))
nu = n - 2
  \end{codigo}
  \begin{block}{}
    \resalta{Incertidumbres de los parámetros:}
    \begin{equation}
      u(a)=\frac{1}{\sqrt{\sum_i w_i}\,\sqrt{\langle x^{2}\rangle-\langle x\rangle^{2}}},
      \qquad
      u(b)=u(a)\sqrt{\langle x^{2}\rangle}.
    \end{equation}
  \end{block}
  \small
  \begin{itemize}
    \item $b=-0{,}1014$~m/s: la recta pasa por el centro de gravedad ponderado.
    \item $u(a)=7{,}67$~m/s, $u(b)=0{,}0062$~m/s y $\nu=6-2=4$.
  \end{itemize}
\end{frame}

\begin{frame}[fragile]{Bondad del ajuste: $\chi^{2}$}
  \begin{codigo}[listing options={style=python, basicstyle=\ttfamily\scriptsize\color{codtexto}}]
chi2 = float((w * (y - a * x - b) ** 2).sum())
factor = max(1.0, math.sqrt(chi2 / nu)) if nu > 0 else 1.0
  \end{codigo}
  \begin{block}{}
    \begin{equation}
      \chi^{2}=\sum_i\left(\frac{y_i-ax_i-b}{u_i}\right)^{2}.
    \end{equation}
  \end{block}
  \small
  \begin{itemize}
    \item Si las barras de error están bien estimadas, $\chi^{2}\approx\nu$.
    \item En el caso~1, $\chi^{2}=69{,}3$ frente a $\nu=4$: la dispersión supera lo
      que indican las barras y $u(a)=7{,}67$~m/s es demasiado optimista.
    \item Se corrige con $\sqrt{69{,}3/4}=4{,}16$; \cod{max(1.0, ...)} impide que la
      incertidumbre disminuya.
  \end{itemize}
\end{frame}

\begin{frame}[fragile]{La $t$ de Student y el resultado}
  \begin{codigo}[listing options={style=python, basicstyle=\ttfamily\scriptsize\color{codtexto}}]
t95 = float(stats.t.ppf(0.975, nu)) if nu > 0 else float("nan")
return dict(a=a, b=b, ua=ua * factor, Ua=t95 * ua * factor, ...)
  \end{codigo}
  \cod{stats.t.ppf(0.975, 4)} deja un $2{,}5$\,\% en cada cola, es decir, un
  intervalo central del 95\,\%: $t_{95,4}=2{,}78$, mayor que el $1{,}96$ de la
  gaussiana, porque con pocos puntos la propia incertidumbre está mal conocida.
  \begin{center}\small
  \begin{tabular}{@{}lll@{}}
    \toprule
    Clave & Contenido & Caso 1\\
    \midrule
    \texttt{a} & pendiente & $244{,}657$~m/s\\
    \texttt{ua} & $u(a)$ corregida & $7{,}67\times4{,}16=31{,}9$~m/s\\
    \texttt{Ua} & $t_{95,4}\,u(a)$ & $2{,}78\times31{,}9=88{,}6$~m/s\\
    \texttt{b}, \texttt{Ub} & ordenada & $-0{,}101\pm0{,}071$~m/s\\
    \bottomrule
  \end{tabular}
  \end{center}
\end{frame}

% =========================================================================
\section{Resultado}
% =========================================================================
\begin{frame}{Del código al resultado}
  \begin{center}
    \scriptsize
    \begin{tikzpicture}
  \begin{axis}[datos, scale only axis, name=principal, width=8.2cm, height=2.7cm,
      xmin=-1, xmax=2.6, ymin=-0.45, ymax=0.45, clip=true,
      xlabel={}, xticklabels={}, ylabel={$v$ (\si{\meter\per\second})},
      legend pos=north west, ytick={-0.4,-0.2,0,0.2,0.4},
      yticklabel style={/pgf/number format/fixed, /pgf/number format/precision=1}]
    \addplot[teorica, domain=-1:2.6, samples=2] {0.34485*x};
    \addlegendentry{$v=\vst\,x$ (valor teórico)}
    \addplot[ajuste, domain=-1:2.6, samples=2] {0.2446553*x+(-0.1014381)};
    \addlegendentry{ajuste ponderado}
    \addplot[puntoA] coordinates {(2.17874,0.35317) +- (0.12257,0.00258) (1.22569,0.22417) +- (0.04482,0.00092) (0.68085,0.12717) +- (0.10117,0.00042)};
    \addlegendentry{acercamiento}
    \addplot[puntoB] coordinates {(-0.41471,-0.24200) +- (0.06189,0.00083) (-0.08620,-0.12833) +- (0.03217,0.00068)};
    \addlegendentry{alejamiento}
    \draw[-{Stealth[length=5pt]}, azulacero, line width=0.8pt]
      (axis cs:-0.55,-0.3825) -- (axis cs:-0.97,-0.3825);
  \end{axis}
  \begin{axis}[datos, at={(principal.south west)}, anchor=north west, yshift=-4pt,
      scale only axis, width=8.2cm, height=0.8cm, xmin=-1, xmax=2.6, ymin=-5, ymax=8.5,
      xlabel={$10^{3}\,x=10^{3}\,(1-\fE/\fR)$}, ylabel={$r_i/u_i$},
      ytick={-3,0,3,6}]
    \addplot[grispizarra, line width=0.5pt, densely dashed, domain=-1:2.6, samples=2] {0};
    \addplot[only marks, mark=*, mark size=1.7pt, azulprof,
      error bars/.cd, y dir=both, y explicit,
      error bar style={line width=0.5pt, azulprof}, error mark options={rotate=90, mark size=1.8pt, azulprof}]
      coordinates {(2.17874,-2.6060) +- (0,1) (1.22569,2.3389) +- (0,1) (-0.41471,-2.5786) +- (0,1) (0.68085,2.5059) +- (0,1) (-0.08620,-0.7351) +- (0,1)};
    \draw[-{Stealth[length=5pt]}, azulacero, line width=0.8pt]
      (axis cs:-0.55,6.602) -- (axis cs:-0.97,6.602);
      \end{axis}
\end{tikzpicture}

  \end{center}
  \begin{block}{}
    \centering
    $\vs=(2{,}4\pm0{,}9)\times10^{2}$~m/s, \qquad $b=(-0{,}10\pm0{,}07)$~m/s\\[1pt]
    {\footnotesize Incertidumbres ampliadas al 95\,\%, con $\nu=4$.}
  \end{block}
\end{frame}

\begin{frame}{Resumen}
  \begin{block}{}
    \resalta{Redondeo:}
    \begin{enumerate}
      \item el logaritmo y \texttt{floor} localizan la primera cifra;
      \item dividir y truncar indica cuál es;
      \item si es un 1 se conservan dos cifras; si no, una;
      \item valor e incertidumbre se cortan en la misma posición.
    \end{enumerate}
  \end{block}
  \begin{block}{}
    \resalta{Ajuste ponderado:}
    \begin{enumerate}
      \item cada punto pesa $1/u_i^{2}$, con la varianza efectiva;
      \item se itera hasta que la pendiente no cambia;
      \item $b$, $u(a)$ y $u(b)$ se calculan con medias ponderadas;
      \item $\chi^{2}/\nu$ indica si hay que corregir las incertidumbres;
      \item la $t$ de Student da el intervalo al 95\,\%.
    \end{enumerate}
  \end{block}
\end{frame}

\end{document}
